\documentclass[11pt,reqno]{amsart}
\usepackage[margin=1.15in]{geometry}
\usepackage{amsmath,amssymb,amsthm,mathtools}
\usepackage{booktabs}
\usepackage[colorlinks=true,linkcolor=blue,citecolor=blue,urlcolor=blue]{hyperref}
\usepackage{microtype}
% allow inter-paragraph stretch so page breaks are not underfull
\raggedbottom
% let hyperref break long URLs at punctuation instead of running into the margin
\usepackage{xurl}

\theoremstyle{plain}
\newtheorem{theorem}{Theorem}
\newtheorem{proposition}[theorem]{Proposition}
\newtheorem{corollary}[theorem]{Corollary}
\newtheorem{lemma}[theorem]{Lemma}
\newtheorem{conjecture}[theorem]{Conjecture}
\theoremstyle{definition}
\newtheorem{definition}[theorem]{Definition}
\newtheorem{remark}[theorem]{Remark}
\newtheorem{observation}[theorem]{Observation}

\DeclareMathOperator{\sqf}{sqf}
\newcommand{\Art}{\mathrm{A}}
\newcommand{\Z}{\mathbb{Z}}
\newcommand{\Q}{\mathbb{Q}}

\begin{document}

\title[Cross-base correlations of Artin primes]
{Cross-base correlations of Artin primes:\\ entanglement, exclusion, and the
structure of $p-1$}

\author{Josh Bald}
\address{Ontario, Canada}
\email{jpbald93@gmail.com}

\subjclass[2020]{Primary 11A07; Secondary 11N05, 11N13, 11Y16, 11Y60}
\keywords{Artin's conjecture, primitive roots, simultaneous primitive roots,
entanglement, quadratic characters, Kummer theory, computational number theory}

\date{\today}

\begin{abstract}
Say a prime $p$ is \emph{Artin base $a$} if $a$ is a primitive root modulo
$p$. This paper studies the joint behaviour of Artin status \emph{across
bases} at the same prime. We first record a deterministic exclusion law: for
\emph{any} two bases $a, b$, writing $d = \sqf(ab)$ for the squarefree part of
$ab$, no prime $p$ with $\bigl(\tfrac{d}{p}\bigr) = -1$ is simultaneously
Artin for $a$ and for $b$. That condition is a union of residue classes
modulo $4d$, so whenever $d > 1$ --- equivalently, whenever $a$ and $b$ do not
have the same squarefree part --- half of all primes are barred outright;
the proof is two lines from the multiplicativity of quadratic characters. A
corollary is the triple form of the law: if $\sqf(c) = \sqf(ab)$ then no odd
prime is simultaneously Artin for $a$, $b$ and $c$ --- for instance, no odd
prime has $2$, $5$ and $10$ all as primitive roots. Over all $66$ pairs of the base set
below and all $50{,}847{,}524$ primes below $10^9$, the number of primes
violating the pair law is zero, as it must be.
We then measure, for all $66$ pairs from the twelve bases
$\{2,3,5,6,7,10,11,13,15,17,21,29\}$ and all $50{,}847{,}532$ primes up to
$10^9$, the same-prime correlation $\phi(a,b)$ between the indicator variables
of Artin status. Three empirical findings emerge. First, the correlation is
positive for $64$ of the $66$ pairs (mean $\phi = 0.0352$, individual $\phi/\sigma_0$ up to $476$ against the benchmark scale $\sigma_0 = N^{-1/2}$ --- which reflects the $5 \times 10^7$ primes enumerated rather than a large effect; the correlations explain about $0.1\%$ of variance),
despite the exclusion law: the common cause $p-1$ dominates. The only negative
pairs are $(2,10)$ and $(3,15)$, the two for which the barring character has
conductor $5$, the smallest occurring. Second, conditioning on the number of prime factors
$\omega(p-1)$ removes only $55\%$ of the mean correlation, but conditioning on
a finer signature of $p - 1$ --- the $2$-adic valuation, the divisibility of
$p-1$ by each prime up to $13$, and the number of larger factors --- removes
$95\%$: the cross-base correlation is, almost entirely, a reflection of the
fact that all bases see the same factorisation of $p - 1$. Third, the residual
that survives fine conditioning is concentrated on the triple-completed pairs (mean absolute residual $0.018$ within those pairs, or
$0.008$ excluding one large positive outlier, versus $0.001$ for
non-completed pairs), and is localised in the quadratic-character
channel, which is where the exclusion law of Theorem~\ref{thm:pair} acts;
we read this as the law's statistical trace, visible once the common cause is
conditioned away, while noting that a support restriction of this kind does
not by itself determine the size or sign of a residual. Finally we evaluate the joint
densities predicted by the Matthews--Moree--Stevenhagen Kummer-degree heuristic for all
$66$ pairs: with no fitted parameters they match the measured densities to a
mean relative error of $0.012\%$, and they reproduce the sign reversal on the two negatively correlated pairs. Both the pointwise obstruction and the Kummer-degree
heuristic are governed by the same quadratic class $\sqf(ab)$, which appears
in the first as a congruence barring half the primes and in the second as a
member of the group of quadratic classes that can reduce the degree. We do not
claim these are the same phenomenon: the product class belongs to that group
for every pair, whereas it lowers a degree only when its conductor divides the
relevant cyclotomic level, and we exhibit pairs separating the two.
\end{abstract}

\maketitle

\section{Introduction}
\label{sec:intro}

Artin's conjecture on primitive roots concerns one base at a time: for
non-square $a \neq -1$, the set
\[
  \Art_a = \{\, p \text{ prime} : a \text{ is a primitive root mod } p \,\}
\]
should have density given by a rational multiple of Artin's constant, as
proved under GRH by Hooley~\cite{Hooley1967}. The joint distribution for
\emph{several} bases at the same prime was treated by
Matthews~\cite{Matthews1976}, who computed (under GRH) the density of primes
for which several prescribed integers are simultaneously primitive
roots.\footnote{We record one limitation of our own scholarship:
\cite{Matthews1976} was verified from the publisher's bibliographic record and
from independent restatements of its main theorem, not from the original text,
which the author was unable to obtain. The same applies to~\cite{Lenstra1977},
which is likewise behind a paywall the author could not clear; the
characterisation of its scope given below follows secondary accounts rather
than the original. Every other reference was checked
against the published record, except the companion
preprints~\cite{BaldI, BaldII}, which are unrefereed.}
Lenstra's Galois-theoretic framework~\cite{Lenstra1977} covers a single
finitely generated group $W$ with a prescribed index condition, and Lenstra explicitly lists the simultaneous problem --- ``the set of primes $q$ for which
both $2$ and $3$ are primitive roots'', citing Matthews --- among
the sets to which his results do not directly apply. The systematic Kummer-degree treatment of the simultaneous densities is given by Moree--Stevenhagen~\cite{MoreeStevenhagen2014}, whose character-sum method
recovers Matthews' theorem; see also Moree's survey~\cite{Moree2012}. The qualitative message of
that theory is that the events $p \in \Art_a$ and $p \in \Art_b$ are
\emph{not} independent, and that the deviation from independence is governed
by the interaction of the Kummer--Galois extensions attached to $a$ and $b$
--- in modern language, by \emph{entanglement} between them.

That framework remains actively developed.
J\"arviniemi--Perucca--Sgobba~\cite{JarviniemiPeruccaSgobba2025} give a unified
treatment of Artin-type problems for finitely generated groups in number
fields; Kimmel~\cite{Kimmel2024} treats number-field and matrix variants (all results there conditional on GRH);
Klurman--Shparlinski--Ter\"av\"ainen~\cite{KlurmanShparlinskiTeravainen2025}
obtain results on average over the base; Sgobba~\cite{Sgobba2025} establishes
unconditional results in cases where GRH can be circumvented; and
Goldmakher--Martin--P\'eringuey~\cite{GoldmakherMartinPeringuey2025} propose
refined conjectures on the distribution of the multiplicative order, which
bear on the structure of $p-1$ that plays the central role in
Section~\ref{sec:decomposition} below.

The density theory for several bases at a fixed prime is well developed:
\cite{JarviniemiPeruccaSgobba2025} treats, under GRH, the joint index map
$\mathfrak{p} \mapsto (\mathrm{Ind}_{\mathfrak{p}}(W_1), \ldots,
\mathrm{Ind}_{\mathfrak{p}}(W_n))$ in full generality, of which the
simultaneous primitive root problem over $\Q$ is the special case
$H = \{(1,\ldots,1)\}$. What appears not to have been done is the empirical
side: measuring these joint densities and the resulting same-prime
correlations at scale, and decomposing them by conditioning on the
factorisation of $p-1$. That is our subject; we
take~\cite{Matthews1976, MoreeStevenhagen2014} as the source of the density
predictions tested in Section~\ref{sec:model}. As those predictions are
conditional on GRH, so is every comparison we make with them; our measurements
themselves are of course unconditional.

This paper is the third in a series studying correlations of Artin status
computationally at the scale of $10^9$~\cite{BaldI, BaldII}. The first
paper~\cite{BaldI} measured the correlation between the Artin statuses (base $10$) of
\emph{consecutive} primes and decomposed it into residue channels; the second
gives a classification (companion preprint, not yet refereed) of deterministic exclusion laws for consecutive Artin
primes in arbitrary bases and identified the conductor of the quadratic
character $\chi_a$ as the parameter governing the statistical correlation.
Here we complete the picture along the remaining axis: \emph{fix the prime,
vary the base}.

Two questions organise the paper.

\emph{(Q1) What is deterministically forbidden?} For consecutive primes, the
answer given in~\cite{BaldII} was a set of gap classes on which the
quadratic character is forced to flip. For a single prime and several bases,
the answer is a congruence condition on $p$ that bars joint Artin status on
half of all primes, for every pair of bases not sharing a squarefree part:

\begin{theorem}[Same-prime pair exclusion; see Theorem~\ref{thm:pair}]
Let $a, b$ be positive integers and set $d = \sqf(ab)$. Then no odd prime $p$
with $\bigl(\tfrac{d}{p}\bigr) = -1$ is simultaneously Artin base $a$ and
Artin base $b$.
\end{theorem}

By quadratic reciprocity the excluded set is an explicit union of residue
classes modulo $4d$, and whenever $d > 1$ it has density exactly $1/2$ among
the odd primes. So for any two bases with $\sqf(a) \ne \sqf(b)$ --- with no
requirement that they sit in a special configuration, and with no third base in
sight --- half the primes are deterministically forbidden from being Artin for
both. (If $\sqf(a) = \sqf(b)$ then $d = 1$, the character is principal and the
statement is vacuous; see Section~\ref{sec:theorem}.) For example $2$ and $6$
are never simultaneously primitive roots for $p \equiv 5, 7 \pmod{12}$; $2$
and $10$ never for $p \equiv 2, 3 \pmod 5$.

Specialising the barring character to a third base of the set gives the triple
form, which was our original route to the law:

\begin{corollary}[Triple exclusion; see Theorem~\ref{thm:triple}]
Let $a, b, c$ be positive integers with $\sqf(c) = \sqf(ab)$. Then no odd prime
is simultaneously Artin base $a$, Artin base $b$, and Artin base $c$.
\end{corollary}

For example, no odd prime has all of $2, 5, 10$ as primitive roots; none has
all of $3, 5, 15$; and, since $\sqf$ is all that matters, none has all of
$2, 5, 40$. (The restriction to odd $p$ is essential and not cosmetic: see the
sharpness remark below.) Like the gap exclusion law of~\cite{BaldII}, both statements are
elementary: they need only the multiplicativity of the quadratic character and
the fact that a quadratic residue is not a primitive root. The underlying criterion
is moreover classical in the density theory. Moree and
Stevenhagen~\cite[\S5, (5.2)]{MoreeStevenhagen2014} characterise when the
local quadratic admissibility set $S_2$ is empty --- equivalently, when the
naive density vanishes --- by the condition that some odd-size subset $I$ has
$\prod_{i \in I} a_i$ a square; the criterion is originally
Matthews'~\cite{Matthews1976}. We emphasise what that statement is and is
not. Emptiness of $S_2$ gives an unconditional obstruction at every odd prime,
and Theorem~\ref{thm:triple} is the case $n = 3$, $I = \{a,b,c\}$ of exactly
that obstruction. The converse --- that the set of primes is infinite, or has
positive density, whenever no such relation holds --- is a statement about
densities and is conditional on GRH in the cited framework, as the authors are
careful to distinguish. We claim no priority for the obstruction itself. What
we add is an elementary pointwise formulation valid at every odd prime with no
hypothesis of any kind, its extension to square bases and to $1$, a
machine-checked proof, and --- the point of emphasis here --- the observation
that the third base is inessential: the obstruction is a congruence, of
density $1/2$, attached to the single character $\chi_{\sqf(ab)}$. That character is the object the rest of the paper measures.

\emph{(Q2) What does the joint distribution actually look like, and what
explains it?} Here the data speak. We compute the exact $2 \times 2$
contingency table of $\bigl(\mathbf{1}[p \in \Art_a],\,
\mathbf{1}[p \in \Art_b]\bigr)$ over all $50{,}847{,}532$ primes
$p \le 10^9$, for all $66$ pairs from the base set
\[
  \{2, 3, 5, 6, 7, 10, 11, 13, 15, 17, 21, 29\},
\]
chosen to contain many multiplicative triples ($15$ of the $66$ pairs lie in
a triple $\sqf(ab) = \sqf(c)$ within the set) alongside pairs the set does not complete. The findings, developed in
Sections~\ref{sec:measurement}--\ref{sec:entanglement}:

\begin{enumerate}
  \item \emph{The exclusion law is visible exactly.} For every one of the $66$
    pairs, the number of primes $p \le 10^9$ that are Artin for both $a$ and
    $b$ while $\chi_a(p) \ne \chi_b(p)$ is $0$, as Theorem~\ref{thm:pair}
    requires; on the complementary half of the primes the joint density is
    $0.299$. The deterministic law is thus not merely proved but
    exhibited, on $50{,}847{,}524$ primes and $3.4 \times 10^9$ pair-prime
    incidences.
  \item \emph{Positivity, nonetheless.} $\phi(a,b) > 0$ for $64$ of $66$
    pairs, with mean $0.0352$, the largest exceeding the i.i.d.\ benchmark scale
$\sigma_0$ by a factor of $476$: knowing that $a$ is a
    primitive root mod $p$ makes it \emph{more} likely that $b$ is, even though
    half of all primes are barred from having both. The common cause $p-1$
    outweighs the exclusion. The largest correlations
    all involve base $5$ (e.g.\ $\phi(5,13) = 0.0667$:
    $P(\Art_{13} \mid \Art_5) = 0.416$ against
    $P(\Art_{13} \mid \lnot\Art_5) = 0.350$).
  \item \emph{Where the exclusion wins.} The only two negative
    correlations are $\phi(2,10) = -0.0302$ and $\phi(3,15) = -0.0303$. Both are triple-completed, but triple-completion alone is not the discriminant: fifteen of the $66$ pairs are triple-completed and thirteen of those are positively correlated. What singles out these two is the conductor of
    the barring character $\chi_{\sqf(ab)}$: both have $\sqf(ab) = 5$, of
    conductor $5$ --- the smallest conductor occurring among the triple-completed pairs, and the only odd prime one. A plausible reading, recorded as an observation in Section~\ref{sec:measurement} rather than proved, is that a small conductor is what lets the exclusion of Theorem~\ref{thm:pair} survive the averaging over $p-1$ instead of being diluted by it.
  \item \emph{The positive correlation is $p-1$ structure, not entanglement.}
    Conditioning on $\omega(p-1)$ removes $55\%$ of the mean correlation;
    conditioning on a fine signature of $p-1$ --- namely
    $\bigl(\min(v_2(p-1), 4),\ \text{the set of primes } q \le 13 \text{
    dividing } p-1,\ \min(\omega_{>13}(p-1), 3)\bigr)$ --- removes $95\%$,
    leaving a mean residual of $0.0019$. All bases are tested against the
    same integer $p - 1$; a prime whose $p-1$ is ``primitive-root friendly''
    (few small factors, small $2$-part) is friendlier for \emph{every} base
    simultaneously. This common-cause structure, not any interaction between
    the bases, is the bulk of the correlation.
  \item \emph{What survives is entanglement, and it is exactly localised.}
    After fine conditioning, the residual correlation among \emph{triple-completed} pairs (mean absolute residual $0.0179$, or $0.0077$
    excluding the single large positive outlier $(2,6)$, with individual values
    reaching $-0.034$ for $(2,10)$ and $-0.026$ for $(5,15)$) is eighteen times that of the remaining pairs ($0.0010$), or eightfold with $(2,6)$ excluded;
    and additionally conditioning on the quadratic-residue statuses
    $(\chi_a(p), \chi_b(p))$ flattens triple-completed and remaining pairs to a common $+0.005$. The residual is thus localised, empirically, in the quadratic-character channel $\chi_a \chi_b = \chi_{\sqf(ab)}$ --- the
    statistical shadow of Theorem~\ref{thm:pair}.
  \item \emph{Theory reproduces all of it.} Evaluating the Matthews--Moree--Stevenhagen
    Kummer-degree heuristic~\eqref{eq:model} for each pair
    (Section~\ref{sec:model}) predicts the measured joint densities to a mean
    relative error of $0.012\%$ over the $66$ pairs, with no fitted
    parameters, including the negative $\phi$ of the two negatively correlated pairs. The degree collapse in the model, where it occurs, is driven by the same quadratic class $\sqf(ab)$ that bars half the primes in Theorem~\ref{thm:pair}; the class is present for every pair, while the collapse is not.
\end{enumerate}

The theme of the series persists: at every level, the joint behaviour of
Artin primes decomposes into a small deterministic core (here, the pair
exclusion law) plus a statistical component with an identifiable arithmetic
carrier (here, the shared factorisation of $p-1$, with entanglement as a
measurable second-order correction). What this paper adds to the pattern is that one quadratic class, $\sqf(ab)$, is common to both descriptions: it bars half the primes pointwise in Theorem~\ref{thm:pair}, and it is the class whose conductor, when it divides the relevant cyclotomic level, produces the degree collapse that bends the predicted densities.

\section{The same-prime exclusion laws}
\label{sec:theorem}

Throughout, for a positive integer $a$ we write $d_a = \sqf(a)$
for its squarefree part and $\chi_a = \bigl(\tfrac{d_a}{\cdot}\bigr)$ for the
quadratic character attached to $\Q(\sqrt{a})$; thus $\chi_a(p) = +1$ exactly
when $a$ is a quadratic residue modulo $p$ (for $p \nmid 2a$). We allow $a$ to
be a square, in which case $d_a = 1$ and $\chi_a \equiv +1$ is the principal
character; this convention costs nothing and lets the statements below dispense
with a non-squareness hypothesis. We have
\begin{equation}
\label{eq:obstruction}
  p \in \Art_a \implies \chi_a(p) = -1,
\end{equation}
since a quadratic residue generates a subgroup of index at least $2$ in
$(\Z/p\Z)^\times$.

Everything in this section follows from one line: the quadratic character is
completely multiplicative, so for any positive integers $a, b$ and any odd
prime $p \nmid ab$,
\begin{equation}
\label{eq:mult}
  \chi_a(p)\,\chi_b(p)
    = \Bigl(\tfrac{d_a d_b}{p}\Bigr)
    = \Bigl(\tfrac{\sqf(ab)}{p}\Bigr)
    = \chi_{\sqf(ab)}(p),
\end{equation}
since $d_a d_b$ and $\sqf(ab)$ differ by a square. Combining
\eqref{eq:mult} with the obstruction~\eqref{eq:obstruction} gives the
following.

\begin{theorem}[Pair exclusion]
\label{thm:pair}
Let $a, b$ be positive integers and set $d = \sqf(ab)$. Then for every odd
prime $p$ with $\bigl(\tfrac{d}{p}\bigr) = -1$, the prime $p$ is not
simultaneously Artin base $a$ and Artin base $b$.
\end{theorem}

\begin{proof}
If $p \mid ab$ then one of the two bases is $\equiv 0 \pmod p$ and so is not a
primitive root; assume therefore $p \nmid ab$. By~\eqref{eq:mult},
$\chi_a(p)\chi_b(p) = \bigl(\tfrac{d}{p}\bigr) = -1$, so $\chi_a(p)$ and
$\chi_b(p)$ have opposite signs and in particular one of them equals $+1$;
by~\eqref{eq:obstruction} the corresponding base is not a primitive root.
\end{proof}

The condition $\bigl(\tfrac{d}{p}\bigr) = -1$ is a condition on residue
classes. For $d > 1$ the character $\bigl(\tfrac{d}{\cdot}\bigr)$ is
non-principal, of fundamental discriminant $d$ if $d \equiv 1 \pmod 4$ and
$4d$ otherwise; either way the discriminant divides $4d$, so the excluded set
is an explicit union of reduced classes modulo $4d$, and exactly half of those
classes carry character $-1$. By the prime number theorem in arithmetic
progressions the excluded set then has density exactly $1/2$ among the odd
primes. So for any two bases with $d > 1$ --- no third base required, no
special configuration --- half of all primes are deterministically barred from
being Artin for both, and Theorem~\ref{thm:pair} is a \emph{congruence}
exclusion law.

The excluded case $d = 1$ is exactly $\sqf(a) = \sqf(b)$, i.e.\ $ab$ a perfect
square, i.e.\ $\chi_a = \chi_b$: then $\chi_a \chi_b = \chi_a^2 \equiv +1$,
nothing is ever barred, and Theorem~\ref{thm:pair} is vacuous rather than
false. (For instance $a = 2$, $b = 8$: the two bases generate the same
quadratic field, and indeed $2$ and $8$ are simultaneously primitive roots for
$p = 3, 5, 11, 29, 53, \ldots$) Such pairs are precisely the ones our base set
avoids: all twelve bases are squarefree and distinct, so $d > 1$ for all $66$
pairs. Writing the classes out in the four cases we return to below:
\[
  \begin{aligned}
    (2,6),\; d = 3&\!: \ p \equiv 5, 7 \!\!\pmod{12}; &\qquad
    (3,6) \text{ and } (5,10),\; d = 2&\!: \ p \equiv 3, 5 \!\!\pmod 8; \\
    (2,10),\; d = 5&\!: \ p \equiv 2, 3 \!\!\pmod 5 .
  \end{aligned}
\]
We verified each over all primes below $6 \times 10^4$ with no exceptions, and
Section~\ref{sec:measurement} verifies the theorem itself against all
$50{,}847{,}524$ primes below $10^9$ for all $66$ pairs of our base set
simultaneously. Theorems~\ref{thm:pair} and~\ref{thm:triple} are also
machine-checked in Lean~4 against Mathlib, including the
step~\eqref{eq:obstruction} and the $(5,10)$ instance above, whose barring
condition is discharged from $p \bmod 8$ --- so that instance is verified with
no character hypothesis assumed at all. Every checked statement depends on no
axiom beyond \texttt{propext}, \texttt{Classical.choice} and
\texttt{Quot.sound}. We state the limit of this plainly: the general
translation of $\bigl(\tfrac{d}{p}\bigr) = -1$ into residue classes
modulo $4d$ is \emph{not} formalised generically, but carried out by hand per
instance. See \texttt{LEAN\_NOTE.md} in the accompanying package.

The triple form of the law, which was our original route to it, is the special
case in which the barring character is supplied by a third base of the set.

\begin{theorem}[Triple exclusion]
\label{thm:triple}
Let $a, b, c$ be positive integers with $\sqf(c) = \sqf(ab)$ and let $p$ be an
odd prime. Then $p$ is not simultaneously Artin base $a$, Artin base $b$, and
Artin base $c$.
\end{theorem}

\begin{proof}
If $p$ is Artin base $c$ then $\chi_c(p) = -1$ by~\eqref{eq:obstruction}, and
$\chi_c = \chi_d$ with $d = \sqf(ab) = \sqf(c)$; now apply
Theorem~\ref{thm:pair}.
\end{proof}

\begin{corollary}
\label{cor:twoofthree}
Within any multiplicative triple --- for instance $(2, 5, 10)$, $(3, 5, 15)$,
$(3, 7, 21)$, $(2, 3, 6)$ --- an \emph{odd} prime can be Artin for at most two
of the three bases. The restriction to odd $p$ is necessary and not cosmetic:
every odd integer is vacuously a primitive root modulo $2$, so all three
members of $(3,5,15)$ are Artin base $2$.
\end{corollary}

The gain in passing from Theorem~\ref{thm:triple} to Theorem~\ref{thm:pair} is
not cosmetic, though it is a gain in usable content rather than in logical
strength: both are pointwise statements at every odd prime. The difference is
in what one must know to apply them. To invoke
Theorem~\ref{thm:triple} one must first identify a third base $c$ with
$\sqf(c) = \sqf(ab)$ \emph{and} establish that $c$ is a primitive root at the
prime in question --- itself an instance of the Artin problem, conjectural in
general and (under GRH) of positive but base-dependent density.
Theorem~\ref{thm:pair} needs neither: the barring condition is a congruence,
decidable from $p$ alone by reciprocity, and for $d > 1$ it holds on a set of
density exactly $1/2$. The triple law is the shadow that the pair law casts when
one insists on reading the barring character off a primitive root instead of off
a residue class.

\begin{remark}[Sharpness of the hypotheses]
In both theorems no hypothesis on the bases beyond positivity is needed, only ``$p$ odd'' (together with the barring condition $(d/p) = -1$ in Theorem~\ref{thm:pair}). Non-squareness of the bases is not
required: if $a$ is a square then $\chi_a(p) = +1$ and $a$ is not a primitive
root, so the conclusion is immediate. The same remark disposes of $c = 1$ in
Theorem~\ref{thm:triple}, which arises whenever $\sqf(a) = \sqf(b)$. The
hypothesis $p \ne 2$ must however be kept, and is not removable:
$(\Z/2\Z)^\times$ is trivial, so every odd integer is vacuously a primitive
root mod $2$ and e.g.\ $(a,b,c) = (3,5,15)$ violates the conclusion there. The accompanying Lean formalisation (Section~\ref{sec:theorem}) confirms
independently that the $a \not\equiv 0$ and $b \not\equiv 0 \pmod p$
assumptions are never used.
\end{remark}

\begin{remark}
The exclusion is sharp: primes Artin for exactly two of $(2,5,10)$ exist in
abundance. The smallest is $p = 3$, where $2$ and $5$ are both primitive roots
(each of order $2 = p-1$) while $10 \equiv 1$ is not; the smallest for which
$10$ itself is Artin is $p = 7$, where $5$ and $10$ are both primitive
roots (each of order $6$) while $2$ has order $3$; likewise $p = 23$, where
$5$ and $10$ are primitive roots and $2$ is a quadratic residue of order
$11$. More systematically, in our data the pair
$(2, 10)$ has $P(\Art_{10} \mid \Art_2) = 0.355$: two-of-three is common,
three-of-three is impossible.
\end{remark}

\begin{remark}
Theorem~\ref{thm:pair} is the same-prime analogue of the gap exclusion law
for consecutive primes proved in~\cite{BaldII}: both are
instances of a parity obstruction in the group generated by the quadratic
characters involved, and in both the excluded set is a union of residue
classes. There the two characters were $\chi_a$ evaluated at $p$
and at $p + g$, and the obstruction was arranged by the shift $g$; here the
characters belong to different bases at the same prime, and the obstruction
is arranged by multiplicative dependence. The Galois-theoretic density
computations of Matthews~\cite{Matthews1976} and Moree--Stevenhagen%
~\cite{MoreeStevenhagen2014} contain this relation as the vanishing criterion
recalled in Section~\ref{sec:intro} (the degree of the compositum of the
relevant Kummer extensions drops); the elementary pointwise form seems worth
recording.
\end{remark}

\section{Data and computation}
\label{sec:data}

By a segmented sieve we enumerated all primes $p \le 10^9$
($50{,}847{,}534$ primes). The programs discard only the smallest primes, for
which the $128$-bit Euler-criterion arithmetic and the signature encoding are
degenerate: $p < 5$ for the correlation and decomposition runs, leaving
$50{,}847{,}532$ primes, and $p < 31$ for the signature-and-character run of
Section~\ref{sec:entanglement}, leaving $50{,}847{,}524$. A prime dividing a
base is \emph{not} discarded; it is simply recorded as non-Artin for that base,
which is correct, since $0$ is not a primitive root. The two counts therefore
differ by exactly the eight primes in $[5, 31)$. For each prime we
factored $p - 1$ once by trial division and computed the Artin indicator for
each of the twelve bases by the standard criterion
($a^{(p-1)/\ell} \not\equiv 1 \bmod p$ for every prime $\ell \mid p-1$, in
$128$-bit arithmetic). One streaming pass accumulates, for each of the $66$
unordered base pairs, the $2 \times 2$ contingency table of joint Artin
status --- globally, and stratified by (i) $\omega(p-1)$ and (ii) the fine
signature
\[
  \mathrm{sig}(p) = \Bigl(\min(v_2(p-1), 4),\;\;
  \{\, q \le 13 : q \mid p-1 \,\},\;\; \min(\omega_{>13}(p-1), 3)\Bigr),
\]
a partition of the primes into at most $4 \cdot 32 \cdot 4 = 512$ cells (for
odd $p$ we always have $v_2(p-1) \ge 1$, so the truncated valuation takes four
values, not five). The
signature records exactly the information about $p-1$ that the Artin
criterion consults for small bases: the $2$-part and the small odd prime
divisors. A separate run, also to $10^9$ ($50{,}847{,}524$ primes), stratified
each pair additionally by the quadratic-residue statuses
$\bigl(\chi_a(p), \chi_b(p)\bigr)$ of the two bases on top of the signature.
The whole computation runs in under two hours on a single core of a
commodity machine ($2$-core Intel Xeon E5-2673~v4, $8$~GB RAM); all counts
are exact.

As validation, the marginal Artin densities agreed with the Hooley densities
to within $6 \times 10^{-5}$ absolute ($0.014\%$ relative) for every base ---
e.g.\ $0.373957$ for base $2$, against Artin's constant
$0.373956$, and $0.393642$ for base $5$, whose density carries the classical
correction factor $C \cdot \tfrac{20}{19}$ for
$d \equiv 1 \pmod 4$ (see~\cite{Hooley1967, Moree2012}).

\section{The cross-base correlation}
\label{sec:measurement}

For each pair write $\phi(a,b)$ for the Pearson correlation of the indicator
variables (the $\phi$ coefficient of the $2\times2$ table). This is an exhaustive census of a fixed
interval for a deliberately chosen finite set of bases, so the coefficients
below are exact and carry no sampling error. We therefore report effect sizes
and their variation with $x$, and avoid the vocabulary of statistical
inference. Where a scale factor is quoted we write
$\sigma_0 = N^{-1/2} \approx 0.00014$ at $N \approx 5.08 \times 10^7$: this is
the standard deviation of $\phi$ under an i.i.d.\ independence benchmark that
does \emph{not} describe these data --- the primes are enumerated, not
sampled --- and ratios $\phi/\sigma_0$ should be read as model-relative
diagnostics indicating that the effects are far larger than that benchmark's
scale, not as tests of any hypothesis. No claim below depends on such a ratio.

One convention should be stated, since it affects the conditional averages in
Sections~\ref{sec:decomposition} and~\ref{sec:entanglement}. Within-cell $\phi$
is undefined when a cell is constant in one of the two indicators (for example
cells with $5 \mid p-1$, in which no prime is Artin base $5$). Such cells carry
zero conditional covariance and are omitted from the pooled average, which
therefore draws on $73\%$ of primes on average for the signature
stratification, and $23\%$ for the finer signature-plus-character
stratification, where three of the four character quadrants are identically
non-Artin. Explicitly, for a base pair $(a,b)$ we report
\[
  \phi \mid \mathrm{sig} \;=\;
  \frac{\sum_{s \in E_{ab}} n_s\,\phi_{ab,s}}{\sum_{s \in E_{ab}} n_s},
\]
where $s$ runs over signature cells, $n_s$ is the number of primes in cell $s$,
and $E_{ab}$ is the set of cells with $n_s \ge 200$ in which \emph{both}
indicators vary. This is a prime-count-weighted descriptive mean of within-cell
correlations, not a partial correlation coefficient and not an additive
component of $\phi$; outer means over the $66$ base pairs are unweighted.
Pooling instead on the covariance scale over all cells, which
requires no such exclusion, gives $95.3\%$ removal, in agreement with
Table~\ref{tab:decomp}.

\begin{table}[t]
\centering
\caption{Same-prime cross-base Artin correlation at $x = 10^9$: summary and
extremes. ``Triple-completed'' means the base set contains a third base $c$ with
$\sqf(c) = \sqf(ab)$. This is a property of the chosen base set, not of the
pair: the two bases of such a pair are multiplicatively independent in the
usual sense, and adding or removing $c$ changes the label but not the pair's
density, character or Kummer field.}
\label{tab:summary}
\begin{tabular}{lrr}
\toprule
 & pairs & mean $\phi$ \\
\midrule
all pairs & $66$ & $+0.0352$ \\
non-completed & $51$ & $+0.0363$ \\
triple-completed & $15$ & $+0.0314$ \\
\midrule
largest: $(5,13)$ & & $+0.0667$ \\
smallest: $(3,15)$ & & $-0.0303$ \\
second smallest: $(2,10)$ & & $-0.0302$ \\
\bottomrule
\end{tabular}
\end{table}

\begin{table}[t]
\centering
\caption{The ten largest $|\phi|$ at $x = 10^9$, with the values after
conditioning on $\omega(p-1)$ and on the fine signature $\mathrm{sig}(p)$
(pair-weighted mean of within-cell $\phi$; cells with fewer than $200$
primes excluded). ``TC''\ marks triple-completed pairs.}
\label{tab:top}
\begin{tabular}{rrrrrc}
\toprule
$a$ & $b$ & $\phi$ & $\phi \mid \omega$ & $\phi \mid \mathrm{sig}$ & TC \\
\midrule
$5$ & $13$ & $+0.0667$ & $+0.0200$ & $+0.0003$ & \\
$5$ & $17$ & $+0.0649$ & $+0.0218$ & $+0.0025$ & \\
$5$ & $29$ & $+0.0631$ & $+0.0246$ & $+0.0007$ & \\
$5$ & $10$ & $+0.0624$ & $+0.0285$ & $-0.0125$ & \checkmark \\
$5$ & $6$  & $+0.0623$ & $+0.0292$ & $+0.0005$ & \\
$5$ & $11$ & $+0.0623$ & $+0.0300$ & $+0.0004$ & \\
$5$ & $15$ & $+0.0622$ & $+0.0297$ & $-0.0255$ & \checkmark \\
$3$ & $5$  & $+0.0622$ & $+0.0284$ & $+0.0015$ & \checkmark \\
$5$ & $7$  & $+0.0622$ & $+0.0297$ & $+0.0002$ & \\
$2$ & $5$  & $+0.0621$ & $+0.0336$ & $+0.0004$ & \checkmark \\
\bottomrule
\end{tabular}
\end{table}

Three observations, in increasing order of depth.

\begin{observation}[Positivity]
\label{obs:positive}
$\phi(a, b) > 0$ for $64$ of the $66$ pairs. Artin statuses of different
bases at the same prime \emph{attract}.
\end{observation}

The sign is the opposite of the consecutive-prime correlation measured in the
first two papers of the series~\cite{BaldI, BaldII} (where $\delta < 0$ uniformly), and the
mechanism is transparent: there the two primes had \emph{different} (and
LOS-anticorrelated) values of $p - 1$; here every base is tested against the
\emph{same} $p - 1$. If $p - 1$ has few small prime factors and a small
$2$-part, the primitive-root condition is easier for every base at once.
Formally, the Artin indicators share the common cause $\mathrm{sig}(p)$.

\begin{observation}[The negative pairs are triple-completed]
\label{obs:negative}
The only pairs with $\phi < 0$ are $(2, 10)$ and $(3, 15)$: in each, the
product of the two bases is, up to a square, the third base of a
multiplicative triple contained in the set ($2 \cdot 10 = 4 \cdot 5$,
$3 \cdot 15 = 9 \cdot 5$). For $(2, 10)$:
$P(\Art_{10} \mid \Art_2) = 0.3551$ versus
$P(\Art_{10} \mid \lnot\Art_2) = 0.3853$.
\end{observation}

The sign reversal is where Theorem~\ref{thm:pair} becomes visible statistically. For every pair, triple-completed or not, the primes split into two halves according to the sign
of $\chi_d(p)$, $d = \sqf(ab)$: on the half with $\chi_d(p) = -1$ the values
$\chi_a(p)$ and $\chi_b(p)$ are forced \emph{opposite}, so joint Artin status
is impossible there, while on the half with $\chi_d(p) = +1$ they are forced
\emph{equal} and joint Artin status is unobstructed. This is visible directly
in the data: over all $66$ pairs and all $50{,}847{,}524$ primes, the number of
primes Artin for both bases with $\chi_a(p) \ne \chi_b(p)$ is exactly zero,
while the conditional joint density on the complementary half is
$0.299$.

What then distinguishes the two negatively correlated pairs is not the
existence of the exclusion --- every one of the $66$ pairs has it, on exactly
half the primes --- but the identity of the barring character $\chi_d$, and in
particular how its conductor interacts with the common cause $p-1$. Among the
fifteen triple-completed pairs, $d = \sqf(ab)$ takes the values
$\{2, 3, 5, 6, 7, 10, 15, 21\}$, and the two pairs with $d = 5$ ---
$(2,10)$ and $(3,15)$, of conductor $5$, the smallest conductor occurring and
the only odd prime one --- are exactly the two with $\phi < 0$. That is the
sharp empirical discriminant: not dependence, which thirteen positively
correlated pairs also enjoy, but $\mathrm{cond}(\sqf(ab)) = 5$. A plausible
reading is that a small conductor makes the barring condition
$\chi_d(p) = -1$ resolve on a coarse residue class that is comparatively well
correlated with the divisibility structure of $p-1$ driving the common cause,
so the exclusion survives the averaging instead of being diluted by it; we do
not have a proof of this and record it as an observation.
Section~\ref{sec:entanglement} tests the reading by conditioning the common
cause away, after which the exclusion becomes visible on six triple-completed pairs rather than two.

\section{Decomposition: the correlation is the structure of \texorpdfstring{$p-1$}{p-1}}
\label{sec:decomposition}

\begin{table}[t]
\centering
\caption{Mean cross-base correlation after conditioning, all $66$ pairs at
$x = 10^9$. ``Removed'' is
$1 - \lvert\overline{\phi_{\mathrm{cond}}}\rvert / \lvert\overline{\phi}\rvert$,
the means taken signed over the $66$ pairs; on the mean-absolute scale the
corresponding figures are $51\%$ and $87\%$, and on the covariance scale
$57\%$ and $95\%$.}
\label{tab:decomp}
\begin{tabular}{lrr}
\toprule
conditioning & mean $\phi$ & removed \\
\midrule
none & $0.0352$ & --- \\
$\omega(p-1)$ & $0.0158$ & $55\%$ \\
$\mathrm{sig}(p)$: $\bigl(v_2,\ \{q \le 13\},\ \omega_{>13}\bigr)$
  & $0.0019$ & $95\%$ \\
\bottomrule
\end{tabular}
\end{table}

Conditioning on $\omega(p-1)$ alone --- the crude size of the factorisation
--- removes barely half of the correlation. This is worth emphasising because
$\omega$ is the statistic through which the consecutive-prime papers routed
the mechanism, and it is genuinely insufficient here: two primes with the same
$\omega(p-1)$ can present very different Artin conditions (it matters
\emph{which} primes divide $p-1$, and with what $2$-adic valuation, not how
many). The fine signature --- $2$-adic valuation capped at $4$, the exact set
of primes $\le 13$ dividing $p - 1$, and the count of larger prime factors
capped at $3$ --- is the information the Artin criterion actually consults
for small bases, and conditioning on it removes $95\%$ of the mean
correlation (Table~\ref{tab:decomp}). That the fine structure of $p-1$, rather
than its coarse size, is what governs Artin behaviour is in the same spirit as
the refined order-distribution conjectures
of~\cite{GoldmakherMartinPeringuey2025}, which pass from the Artin condition
$\mathrm{ord}_p(a) = p-1$ to the distribution of
$\omega\bigl((p-1)/\mathrm{ord}_p(a)\bigr)$ --- finer information about the
index than its mere vanishing, just as our signature is finer information about
$p-1$ than $\omega(p-1)$.

We regard this as the main statistical finding: \emph{cross-base Artin
correlation at a single prime is, to first order, not an interaction between
the bases at all}. It is a common-cause effect of the shared integer $p - 1$.
Any model that treats $\Art_a$ and $\Art_b$ as conditionally independent
given the small-prime structure of $p-1$ reproduces $95\%$ of the observed
association.

\section{The residual and the exclusion law}
\label{sec:entanglement}

The $5\%$ that survives is not uniformly distributed: it is concentrated on the triple-completed pairs. This section argues that it is the statistical trace of Theorem~\ref{thm:pair} --- the
deterministic exclusion that the unconditional correlation of
Section~\ref{sec:measurement} had largely averaged away.

\begin{table}[t]
\centering
\caption{Residual correlation after conditioning on $\mathrm{sig}(p)$ and
after additionally conditioning on the pair of quadratic-residue statuses
$(\chi_a(p), \chi_b(p))$, split by triple-completion. All values at
$x = 10^9$ over $50{,}847{,}524$ primes.}
\label{tab:residual}
\begin{tabular}{lrrr}
\toprule
group & mean $\phi$ & mean $\phi \mid \mathrm{sig}$
      & mean $\phi \mid \mathrm{sig}, \mathrm{QR}$ \\
\midrule
triple-completed ($15$ pairs)   & $+0.0314$ & $+0.0049$ & $+0.0049$ \\
\quad mean \emph{absolute} residual        &           & $\;\;0.0179$ & \\
non-completed ($51$ pairs) & $+0.0363$ & $+0.0010$ & $+0.0052$ \\
\quad mean \emph{absolute} residual        &           & $\;\;0.0010$ & \\
\midrule
all $66$ pairs                            & $+0.0352$ & $+0.0019$ & $+0.0052$ \\
\bottomrule
\end{tabular}
\end{table}

Two facts pin down the mechanism, both measured over the full $10^9$ range.

First, after fine conditioning the triple-completed pairs retain substantial residuals, and every pair with a substantial \emph{negative} residual is
triple-completed. The six pairs below are the only ones with
$\phi \mid \mathrm{sig} < -0.004$, and all six are triple-completed. We state this as
the exact combinatorial fact it is --- six of the six most negative residuals fall in the triple-completed class --- rather than attaching a $p$-value to it: the
$66$ pairs share endpoints and the base set was chosen to contain
multiplicative triples, so the pair labels are neither independent nor
exchangeable and no randomisation model licences such a test.
\[
  \begin{aligned}
    (2,10)&\!:\, -0.0343, &\quad (5,15)&\!:\, -0.0255, &\quad (6,10)&\!:\, -0.0144, \\
    (5,10)&\!:\, -0.0125, &\quad (7,21)&\!:\, -0.0060, &\quad (10,15)&\!:\, -0.0049 .
  \end{aligned}
\]
At $x = 10^9$ these are also the only six negative residuals of any size, but
that sharper statement is not stable with $x$: at $x = 3 \times 10^7$ a further twelve pairs have residuals marginally below
zero, all within $3\sigma_0$ of it in the benchmark scale of
Section~\ref{sec:measurement}.

The residual is not sign-definite, and we record the exceptions explicitly.
Nine of the fifteen triple-completed pairs have \emph{positive} residuals: eight of
them small, and one, $(2,6)$, a large outlier at $+0.1598$ --- the largest
residual of any pair in the collection. Its mechanism is transparent and is the same
entanglement phenomenon seen from the other side: $\sqf(6) = 2 \cdot 3$. The relevant family is narrower than $3 \nmid p-1$ alone, however. If $3 \nmid p-1$ and $v_2(p-1) \ge 2$ then $p \equiv 5 \pmod{12}$, so $(3/p) = -1$ and Theorem~\ref{thm:pair} \emph{forbids} $2$ and $6$ together. Agreement can occur only on the complementary family $v_2(p-1) = 1$, that is $p \equiv 11 \pmod{12}$, where $(3/p) = +1$ and the two quadratic characters coincide; that the remaining odd-prime tests then also very nearly coincide is what we measure, not something the factorisation $\sqf(6) = 2 \cdot 3$ implies on its own, and the signature records whether $3 \mid p-1$ without recording the
resulting near-determinism inside those cells. (The effect is sharpest in the cells where $p-1$ has no small odd prime factor
at all: for $v_2(p-1) = 1$, no prime $q \le 13$ dividing $p-1$ and exactly one
larger factor, the $2 \times 2$ table for $(2,6)$ is exactly diagonal over
$1{,}775{,}686$ primes, giving $\phi = +1$. Across the $44$ cells with
$v_2(p-1) = 1$ and $3 \nmid p-1$ the within-cell $\phi$ ranges from $+0.39$ to
$+0.98$, so the near-determinism is a feature of that family of cells rather
than of a single one.) The triple-completed-group mean absolute residual of $0.0179$ is
substantially driven by this one pair: excluding it, the figure is $0.0077$,
still eightfold the value for the remaining pairs of $0.0010$. Conditioning
additionally on $(\chi_a(p), \chi_b(p))$ reduces $(2,6)$ to $+0.0055$, in line
with every other pair --- so the outlier is explained by the entanglement
account rather than in tension with it.
Note that $(2,10)$ and $(5,15)$ are more negative \emph{after} conditioning
than before, and that $(5,10)$, $(5,15)$, $(6,10)$, $(10,15)$ and $(7,21)$
change sign under conditioning: the common-cause channel had been masking an
underlying repulsion.

Second, additionally conditioning on the pair of quadratic-residue statuses $(\chi_a(p), \chi_b(p))$ removes the negative residuals. We put it this way deliberately: the signed group means do not both fall to zero. The triple-completed mean stays at $+0.0049$ while the mean for the remaining pairs \emph{rises} from $+0.0010$ to $+0.0052$, and the two averages are taken over different eligible populations (about $73\%$ of primes for signature conditioning against $23\%$ here), so they are not components of a single decomposition. Every one of the
six negative values above becomes positive and small
($+0.0031$, $+0.0066$, $+0.0045$, $+0.0065$, $+0.0040$, $+0.0048$
respectively), and the triple-completed and remaining groups become indistinguishable at $+0.0049$ and $+0.0052$ (Table~\ref{tab:residual}).
That common small positive floor of about $+0.005$ is present for all $66$
pairs and is plausibly attributable to residual odd-prime conditions shared through $p-1$ beyond what our signature records. We note that a \emph{quartic} channel is not available as an explanation here: a fourth power is a square, and in the only quadrant contributing defined correlations both bases are quadratic nonresidues, so no further $2$-power residue condition can bind.

Conditioning on the quadratic characters is precisely what removes
character-level entanglement, and that it flattens the triple-completed pairs is consistent with the residual as the statistical trace of the relation
$\chi_a \chi_b = \chi_{\sqf(ab)}$ --- the same object that makes
Theorem~\ref{thm:pair} true.

One part of this can be made exact, though it is a statement about support and not about the residual's size.
Theorem~\ref{thm:pair} asserts that the joint Artin event is empty on the
stratum $\chi_a(p) \ne \chi_b(p)$, and this is what the contingency tables
show: summing the joint-Artin cell over all signature classes and all $66$
pairs, the count of primes $p \le 10^9$ that are Artin for both $a$ and $b$
with $\chi_a(p) \ne \chi_b(p)$ is
\[
  0 \qquad \text{(out of } 50{,}847{,}524 \text{ primes, for every one of the
  } 66 \text{ pairs)},
\]
while on the complementary stratum the joint density is $0.299$. The support restriction is therefore a theorem rather than a conjecture. It does not by itself determine the sign or the magnitude of any residual: a support constraint of this kind is compatible with zero covariance, as the trivial example of two independent fair signs with Artin indicators given by their negative-sign indicators shows. What the residual measures is a quadratic-information component; that this component is the pointwise law's statistical trace is an interpretation the data support, not a consequence of the theorem.

One caveat on the strength of the flattening step, in the other direction.
Within a $(\chi_a, \chi_b)$ stratum the exclusion cannot operate at all: by
criterion~\eqref{eq:obstruction} an Artin prime for $a$ has $\chi_a(p) = -1$, so all the \emph{joint} Artin mass sits in the single quadrant $(-1,-1)$. In each of the other three quadrants at least one of the two indicators is identically zero, so the within-cell correlation is undefined and the conditional covariance vanishes; this is not the same as saying neither indicator varies there. Conditioning on
$(\chi_a, \chi_b)$ therefore \emph{restricts attention to the stratum on which
the constraint of Theorem~\ref{thm:pair} is vacuous}, and the vanishing of the
residual there is in part a structural consequence of that restriction rather
than an independent confirmation. What the step establishes is the
localisation: the residual lives in the character channel and not elsewhere in
the signature. We state it in that weaker form.

The picture that emerges is cleanly layered:
\[
  \underbrace{\text{common cause } p-1}_{\approx 95\%}
  \;+\;
  \underbrace{\text{quadratic entanglement } \chi_a\chi_b = \chi_{\sqf(ab)}}_{\text{residual, localised on triple-completed pairs}}
  \;+\;
  \underbrace{\text{higher-order sharing}}_{\text{small positive floor}}.
\]

\section{Comparison with the predicted joint densities}
\label{sec:model}

The measurements of Sections~\ref{sec:measurement}--\ref{sec:entanglement}
identify the mechanisms empirically. We now check them against theory
quantitatively, by computing the joint densities predicted by the
Matthews--Moree--Stevenhagen framework and comparing with the observed counts pair by
pair.

Write $d_a = \sqf(a)$. The heuristic, in the form originating
with~\cite{Matthews1976} and systematised by the character-sum method
of~\cite{MoreeStevenhagen2014} (with the general framework
in~\cite{JarviniemiPeruccaSgobba2025}), is an inclusion--exclusion over the
degrees of the Kummer extensions that detect the conditions ``$a$ is an
$\ell$-th power residue'':
\begin{equation}
\label{eq:model}
  \mathbf{P}\bigl(p \in \Art_a \cap \Art_b\bigr)
  \;=\; \sum_{m, n \text{ squarefree}}
  \frac{\mu(m)\,\mu(n)}
       {\bigl[\Q(\zeta_L, a^{1/m}, b^{1/n}) : \Q\bigr]},
  \qquad L = \mathrm{lcm}(m,n),
\end{equation}
and the entanglement enters through the degree in the denominator. Generically
\[
  [\Q(\zeta_L, a^{1/m}, b^{1/n}):\Q] = \varphi(L)\,m\,n ,
\]
but the degree drops by the order of the subgroup of quadratic classes
realised inside $\Q(\zeta_L)$, taken from the group
\[
  V(m,n) = \bigl\langle\, d_a \text{ (if $2 \mid m$)},\;
                          d_b \text{ (if $2 \mid n$)} \,\bigr\rangle
  \subseteq \Q^\times/(\Q^\times)^2
\]
whose quadratic field lies inside $\Q(\zeta_L)$, i.e.\ whose conductor divides
$L$. The point of contact with Theorem~\ref{thm:pair} is that when both
$2 \mid m$ and $2 \mid n$ the group $V$ contains not only $d_a$ and $d_b$ but
their product $d_a d_b \sim \sqf(ab)$, so the barring character
$\chi_{\sqf(ab)}$ of Theorem~\ref{thm:pair} is present in the model as a
potential source of degree collapse. Two things must be said precisely here,
because a natural stronger reading is false. First, the product class lies in
$V(m,n)$ for \emph{every} pair of bases whenever both indices are even; its
membership is automatic and has nothing to do with whether a third base of our
set completes a multiplicative triple. Second, membership in $V$ is not
collapse: a class contributes only when its conductor divides $L$. These two
conditions do not coincide, in either direction. For $(7,11)$ --- a pair
outside every triple of our set --- the product class $77$ has conductor $77$,
which divides $L = 154$, while the individual conductors $28$ and $44$ do not,
so the product class does contribute a collapse there. Conversely for $(2,6)$
--- a pair inside the triple $(2,3,6)$ --- the nontrivial classes of $V$ have
conductors $8$, $24$ and $12$, all divisible by $4$, and no squarefree $L$ is,
so no term of~\eqref{eq:model} receives this correction at all. The
triple-completion classification used in our tables is therefore a property of
the chosen base set, not a field-theoretic invariant of the pair, and it is not
equivalent to the occurrence of an extra degree collapse. We evaluate~\eqref{eq:model}
by summing the ramified part exactly over all squarefree $m, n$ supported on
$S = \{2\} \cup \{q : q \mid d_a d_b\}$ and folding the rest into the Euler
product $\prod_{q \notin S}\bigl(1 - \tfrac{2q-1}{q^2(q-1)}\bigr)$,
truncated at $3 \times 10^6$.

As a check on the implementation, the same code specialised to a single base
reproduces the classical Hooley densities: $0.373956$ for bases with
$d \not\equiv 1 \pmod 4$ (Artin's constant), $0.393638$ for $d = 5$,
$0.376368$ for $d = 13$, against measured $0.373957$, $0.393642$, $0.376362$.

\begin{table}[t]
\centering
\caption{Predicted versus measured joint Artin densities, selected pairs at
$x = 10^9$. ``TC''\ marks triple-completed pairs. The final two rows are the two negatively correlated pairs.}
\label{tab:model}
\begin{tabular}{rrrrrrrc}
\toprule
$a$ & $b$ & measured $j$ & model $j$ & rel.\ err.
    & measured $\phi$ & model $\phi$ & TC \\
\midrule
$5$  & $13$ & $0.163939$ & $0.163943$ & $+0.00\%$ & $+0.0667$ & $+0.0667$ & \\
$5$  & $10$ & $0.161962$ & $0.161922$ & $-0.02\%$ & $+0.0624$ & $+0.0623$ & \checkmark \\
$5$  & $15$ & $0.161934$ & $0.161922$ & $-0.01\%$ & $+0.0622$ & $+0.0623$ & \checkmark \\
$3$  & $7$  & $0.149965$ & $0.149971$ & $+0.00\%$ & $+0.0432$ & $+0.0433$ & \checkmark \\
$13$ & $17$ & $0.150249$ & $0.150261$ & $+0.01\%$ & $+0.0383$ & $+0.0383$ & \\
$2$  & $3$  & $0.147361$ & $0.147349$ & $-0.01\%$ & $+0.0321$ & $+0.0321$ & \checkmark \\
$2$  & $6$  & $0.147334$ & $0.147349$ & $+0.01\%$ & $+0.0320$ & $+0.0321$ & \checkmark \\
$7$  & $21$ & $0.144754$ & $0.144728$ & $-0.02\%$ & $+0.0238$ & $+0.0238$ & \checkmark \\
$15$ & $21$ & $0.144725$ & $0.144728$ & $+0.00\%$ & $+0.0237$ & $+0.0238$ & \\
\midrule
$3$  & $15$ & $0.132769$ & $0.132776$ & $+0.01\%$ & $-0.0303$ & $-0.0302$ & \checkmark \\
$2$  & $10$ & $0.132776$ & $0.132776$ & $+0.00\%$ & $-0.0302$ & $-0.0302$ & \checkmark \\
\bottomrule
\end{tabular}
\end{table}

\begin{observation}[The model reproduces the measurements]
\label{obs:model}
Over all $66$ pairs at $x = 10^9$, the mean absolute relative error between
the predicted and measured joint densities is $0.012\%$, with maximum
$0.040\%$ (attained at $(21,29)$). In particular the model predicts the sign reversal: for the negatively correlated pairs $(2,10)$ and $(3,15)$ it gives
$j = 0.132776$ and $\phi = -0.0302$ for both, against measured
$j = 0.132776$, $0.132769$ and $\phi = -0.0302$, $-0.0303$ respectively.
\end{observation}

This is the quantitative counterpart of Sections~\ref{sec:measurement} and~\ref{sec:entanglement}. It is a comparison of unconditional joint densities, so it does not by itself validate any conditional mechanism, but it is a demanding test and the framework passes it. The
$66$-pair pattern --- positive correlation almost everywhere, negative precisely on the two pairs where it is observed to be negative, with magnitudes ordered as observed
--- is not merely \emph{consistent with} the Matthews--Moree--Stevenhagen framework; it is
reproduced by it to within a few parts in $10^4$, with no fitted parameters.
Correspondingly, the same quadratic class $\sqf(ab)$ governs the deterministic law of Theorem~\ref{thm:pair} and enters the heuristic through the group $V(m,n)$. We stop short of calling these one object: the class lies in $V(m,n)$ for every pair once both indices are even, while it lowers a degree only when its conductor divides $L$, and the two conditions come apart in both directions as noted in Section~\ref{sec:model}.

\begin{remark}
\label{rem:eps}
The residual discrepancies are uniformly small (all $66$ pairs below
$0.040\%$) and we have not pursued them; ordinary Hooley-type errors of this
size are expected at $x = 10^9$, since the convergence in Hooley's theorem is
only logarithmic, and our Euler products are truncated at $3 \times 10^6$.
We record one methodological warning, since it cost us a draft. Computing
$\varepsilon(m,n)$ as $2$ raised to the number of classes of $V(m,n)$ with
conductor dividing $L$ --- rather than as the order of the subgroup of $V(m,n)$
realised in $\Q(\zeta_L)$ --- inflates $\varepsilon$ to $8$ on exactly those
pairs with $d_a \equiv d_b \equiv 1 \pmod 4$, where all three of
$d_a, d_b, \sqf(ab)$ have odd conductor. Since $V(m,n)$ has two generators,
$|V(m,n)| \le 4$ and $\varepsilon = 8$ is impossible. The error is invisible in
the single-base specialisation, which is why the check above does not catch it;
it inflates the mean relative error over the $66$ pairs from $0.012\%$ to
$0.037\%$ and concentrates the damage on the ten pairs drawn from
$\{5, 13, 17, 21, 29\}$.
\end{remark}

\section{Discussion}
\label{sec:discussion}

\subsection*{Relation to the first two papers of the series~\cite{BaldI, BaldII}}
The trilogy now covers both axes of the joint distribution of Artin status.
Along the prime axis (consecutive $p_n, p_{n+1}$, one base), correlation is
negative, driven by the Lemke Oliver--Soundararajan bias~\cite{LOS2016}
propagating through the two distinct integers $p_n - 1$ and $p_{n+1} - 1$,
with deterministic exclusion classes in the gap. Along the base axis (one prime, two bases), correlation is positive for $64$ of the $66$ pairs measured, driven by the single shared integer $p - 1$, with deterministic exclusion on an explicit half of all primes for every pair of bases with distinct squarefree parts.
In both cases the deterministic law is a parity statement about quadratic
characters attached to a congruence condition --- there on the gap
$p_{n+1} - p_n$, here on $p$ itself --- and in both cases it is statistically
dominated by a factorisation-of-$p-1$ mechanism: diluted to invisibility in
the global average there, absorbed almost entirely by conditioning here. The
symmetry is closer than we expected when the series began: what looked like a
statement about special configurations of bases (multiplicative triples) turned
out, once the third base was seen to be inessential, to be a statement of the
same shape as the gap law --- a character condition on a residue class of
density $1/2$.

\subsection*{Future work}
(1) A rigorous version of Section~\ref{sec:model}: the comparison there is
between measurement and a heuristic evaluated numerically, and the error terms
--- both the Hooley-type error at $x = 10^9$ and the truncation of our Euler
tail --- are not controlled. Making the agreement into a theorem (conditionally
on GRH, in the manner of~\cite{Matthews1976}) would require uniform error
bounds we do not attempt. (2) The conductor dependence made precise.
Theorem~\ref{thm:pair} bars half the primes for every pair of our base set,
yet only the two pairs with $\mathrm{cond}(\sqf(ab)) = 5$ show a negative
unconditional correlation; Section~\ref{sec:measurement} offers a heuristic reason but no
proof. Quantifying the surviving exclusion as a function of
$\mathrm{cond}(\sqf(ab))$ --- equivalently, computing the covariance between
the barring condition $\chi_{\sqf(ab)}(p) = -1$ and the signature of $p-1$ ---
looks tractable within the Kummer-degree framework and would turn the
empirical discriminant into a formula. A natural test is to enlarge the base
set so that several small odd prime conductors occur among the triple-completed pairs, which our twelve bases do not provide. (3) The mixed axis: $\phi$
between $\Art_a(p_n)$ and $\Art_b(p_{n+1})$ for $a \neq b$ --- the data
pipeline extends immediately, and the prediction from the present findings is a
small negative correlation governed by the two conductors jointly.
(4) Higher-order entanglement: triples of bases conditioned on quadratic data
should expose the cubic layer ($\ell = 3$ Kummer conditions), with
$\Q(\zeta_3)$ playing the role $\Q(i)$ and $\Q(\sqrt{2})$ played here. There is no direct analogue of Theorem~\ref{thm:pair} at $\ell = 3$. The quadratic argument works because $-1$ is the only nonidentity value and the product of two such values is forced to be $+1$; for cubic characters the nonidentity values $\zeta_3, \zeta_3^2$ have products realising all three values, so requiring both bases to be cubic nonresidues imposes no condition on their product. Even a relation $\sqf(c) = \sqf(ab)$ is compatible with the assignment $(\zeta_3, \zeta_3, \zeta_3^2)$. Whether some higher-order Frobenius relation produces an obstruction, and over which population of primes (cubing is a bijection when $p \equiv 2 \bmod 3$), is a question we leave open rather than a stated analogue.

\section*{Acknowledgements}
The author thanks the developers of the open-source tools used in this work
(GCC, Python, \texttt{sympy}, \texttt{matplotlib}).

\subsection*{Use of generative AI}
The author used Anthropic's Claude (models Claude Opus~4.5 and Claude Opus~5)
as an assistant: for drafting and revising exposition, for writing and
debugging the sieve and analysis programs, for exploratory discussion, and for
an adversarial audit of an earlier version of this manuscript, which located
errors that are corrected here. The skeleton of the Lean proof of the
primitive-root/non-residue bridge (\texttt{PrimitiveRootBridge.lean}) was
drafted by a locally hosted open-weights model (Qwen3-32B) and completed by
the author, who supplied the steps the model flagged as uncertain; provenance
is recorded in that file's header. All computations were executed on the
author's own hardware from source code the author has inspected, and numerical
outputs were produced by running deterministic programs, not synthesised by a
language model. The author is solely responsible for the content.

\section*{Declaration of interest}
The author reports there are no competing interests to declare.

\section*{Funding}
This research received no external funding.

\section*{Data availability}
All code and results are publicly available at
\url{https://github.com/jpbald93/consecutive-artin}, including the
single-pass multi-base programs (\texttt{crossbase.c},
\texttt{crossbase\_fine.c}, \texttt{crossbase\_om.c},
\texttt{crossbase\_qr.c}), the analysis scripts, and the raw
contingency-table output at $x = 10^9$ from which every value quoted here is
derived, namely
\begin{center}
\texttt{crossbase\_qr\_1e9.json}, \quad
\texttt{crossbase\_fine\_1e9.json}, \quad
\texttt{crossbase\_om\_1e9.json}
\end{center}
($3.2$~MB in total). Running \texttt{analyze\_qr.py} on the first of these
reproduces Table~\ref{tab:residual} exactly. The full computation reproduces in under two
hours on a single core; the $x = 10^9$ signature-and-character run alone takes
about $16$ minutes on $1$ core of a modern desktop. The Lean sources for the
machine-checked statements are included in the same repository, with build
instructions in \texttt{BUILD.md}.

\begin{thebibliography}{99}

\bibitem{BaldI}
J.~Bald,
\emph{Correlations between primitive root statuses of consecutive primes},
unrefereed preprint (2026), Zenodo,
\href{https://doi.org/10.5281/zenodo.22863946}{doi:10.5281/zenodo.22863946};
code and data at
\url{https://github.com/jpbald93/consecutive-primitive-roots}.

\bibitem{BaldII} J.~Bald, \emph{Quadratic exclusion laws for consecutive Artin primes in arbitrary bases}, unrefereed preprint (2026), Zenodo, \href{https://doi.org/10.5281/zenodo.22865343}{doi:10.5281/zenodo.22865343}; code and data at \url{https://github.com/jpbald93/quadratic-exclusion-laws}.

\bibitem{GoldmakherMartinPeringuey2025}
L.~Goldmakher, G.~Martin, and P.~P\'eringuey,
\emph{Refinements of Artin's primitive root conjecture},
preprint (2025), \texttt{arXiv:2502.19601}.

\bibitem{Hooley1967}
C.~Hooley,
\emph{On Artin's conjecture},
J. Reine Angew. Math. \textbf{225} (1967), 209--220.

\bibitem{JarviniemiPeruccaSgobba2025}
O.~J\"arviniemi, A.~Perucca, and P.~Sgobba,
\emph{Unified treatment of Artin-type problems II},
Res. Number Theory \textbf{11} (2025), no.~42, 19~pp.

\bibitem{Kimmel2024}
N.~Kimmel,
\emph{Artin's primitive root conjecture in number fields and for matrices},
Acta Arith. \textbf{216} (2024), no.~4, 329--348.

\bibitem{KlurmanShparlinskiTeravainen2025}
O.~Klurman, I.~E. Shparlinski, and J.~Ter\"av\"ainen,
\emph{On Artin's conjecture on average and short character sums},
Bull. Lond. Math. Soc. \textbf{57} (2025), 2429--2443.

\bibitem{Lenstra1977}
H.~W. Lenstra, Jr.,
\emph{On Artin's conjecture and Euclid's algorithm in global fields},
Invent. Math. \textbf{42} (1977), 201--224.

\bibitem{LOS2016}
R.~J. Lemke~Oliver and K.~Soundararajan,
\emph{Unexpected biases in the distribution of consecutive primes},
Proc. Natl. Acad. Sci. USA \textbf{113} (2016), no.~31, E4446--E4454.

\bibitem{MoreeStevenhagen2014}
P.~Moree and P.~Stevenhagen,
\emph{Computing higher rank primitive root densities},
Acta Arith. \textbf{163} (2014), no.~1, 15--32.

\bibitem{Matthews1976}
K.~R. Matthews,
\emph{A generalisation of Artin's conjecture for primitive roots},
Acta Arith. \textbf{29} (1976), no.~2, 113--146.

\bibitem{Moree2012}
P.~Moree,
\emph{Artin's primitive root conjecture --- a survey},
Integers \textbf{12A} (2012), Article A13, 100~pp.

\bibitem{Sgobba2025}
P.~Sgobba,
\emph{Unconditional results for Artin-type problems over number fields},
preprint (2025), \texttt{arXiv:2508.08996}.

\end{thebibliography}

\end{document}
